\documentclass[11pt]{amsart}

\usepackage[margin=1.04in]{geometry}
\usepackage{amsmath,amssymb,amsthm}
\usepackage{graphicx}
\usepackage[hidelinks]{hyperref}

\newtheorem{theorem}{Theorem}
\newtheorem{lemma}{Lemma}
\theoremstyle{remark}
\newtheorem{remark}{Remark}

\newcommand{\cb}{\mathrm{cb}}
\newcommand{\op}{\mathrm{op}}
\newcommand{\tr}{\operatorname{tr}}
\newcommand{\id}{\operatorname{Id}}
\newcommand{\C}{\mathbb C}
\newcommand{\Lp}{L^p}
\newcommand{\Sp}{S^p}

\title{A counterexample to the cb-isometry Yeadon-factor conjecture}
\author{Agent lane 14}
\date{Run \texttt{fa\_banach\_001}, 2026-07-18}

\begin{document}
\maketitle

\section*{Status}

This packet gives a candidate counterexample to Conjecture 5.1 and to the
``fundamental'' Conjecture 5.2 of Arhancet \cite{Arhancet2020}.  For every
$1\leq p<\infty$, $p\neq 2$, it constructs finite von Neumann algebras $M$
and $N$ and a positive isometry
\[
        T:\Lp(M)\longrightarrow \Lp(N)
\]
which is completely bounded, but whose Yeadon Jordan homomorphism $J$ is not
completely bounded.  The same example also rules out the structural
decomposition predicted in Conjecture 5.1.

\section*{Source conjectures}

The source paper states the following conjecture in Section 5, page 7.

\begin{center}
\includegraphics[width=0.96\textwidth]{figures/open_problem_crop.png}
\end{center}

In text form, Conjecture 5.1 asks whether, for finite $M$, semifinite $N$ and
$1\leq p<\infty$, $p\neq2$, an isometry $T:\Lp(M)\to\Lp(N)$ is completely
bounded if and only if there is a decomposition $M=M_1\oplus M_2$ and an
integer $n\geq1$ such that $T|_{\Lp(M_1)}$ is a complete isometry and $M_2$ is
$n$-minimal.  The source observes that the forward direction would follow from
Conjecture 5.2: if $T=wbJ$ is a Yeadon factorization and $T$ is completely
bounded, then the normal Jordan $*$-homomorphism $J$ should be completely
bounded.

\section*{Proof Intuition}

The transpose on $S_m^p$ has completely bounded norm
$m^{|1-2/p|}$, although it is an isometry at the scalar $S_m^p$ level.  The
idea is to place both a homomorphic copy and a transposed copy of each matrix
block in the range, but damp the transposed copy on the $m$th block by exactly
$m^{-|1-2/p|}$.  This damping is invisible to scalar isometry after a matching
homomorphic weight is added in $\ell_p$-sum, but it cancels the growth of the
transpose cb norm at all matrix levels.  Thus the isometry $T$ is completely
bounded.

The Yeadon Jordan part does not include the damping weights.  It still
contains the raw transpose on matrix blocks of arbitrarily large dimension,
so it cannot be completely bounded.  Since the same weighted map is not a
complete isometry on any noncommutative matrix block, the proposed
complete-isometry/$n$-minimal decomposition cannot be made: all large blocks
would have to lie in the $n$-minimal side, but their dimensions are unbounded.

\section*{The counterexample}

\begin{theorem}
\label{thm:counterexample}
Let $1\leq p<\infty$ and $p\neq2$.  There exist a finite von Neumann algebra
$M$, a finite hence semifinite von Neumann algebra $N$, and a positive
isometry $T:\Lp(M)\to\Lp(N)$ such that:
\begin{enumerate}
\item $T$ is completely bounded;
\item in the Yeadon factorization $T=wbJ$, the normal Jordan
      $*$-homomorphism $J:M\to N$ is not completely bounded;
\item there is no decomposition $M=M_1\oplus M_2$ with
      $T|_{\Lp(M_1)}$ a complete isometry and $M_2$ $n$-minimal for some
      finite $n$.
\end{enumerate}
Consequently Conjectures 5.1 and 5.2 of \cite{Arhancet2020} are false.
\end{theorem}

\begin{proof}
Put
\[
        \alpha=\left|1-\frac{2}{p}\right|>0.
\]
For $m\geq2$, let $\tr_m$ be the normalized trace on $M_m(\C)$ and let
$\theta_m(x)=x^T$ be matrix transposition.  Define
\[
        M=\prod_{m=2}^{\infty} M_m(\C),\qquad
        N=\prod_{m=2}^{\infty}\bigl(M_m(\C)\oplus M_m(\C)\bigr).
\]
Equip them with the faithful normal finite traces
\[
        \tau_M((x_m))=\sum_{m=2}^{\infty}2^{-m}\tr_m(x_m),
\]
and
\[
        \tau_N((a_m,b_m))
        =\sum_{m=2}^{\infty}2^{-m}\bigl(\tr_m(a_m)+\tr_m(b_m)\bigr).
\]
Both algebras are finite von Neumann algebras; in particular $N$ is
semifinite.  The normalized trace only rescales each finite Schatten block, so
it does not change the completely bounded norms of the block maps used below.

Let
\[
        d_m=m^{-\alpha},\qquad
        c_m=(1-d_m^p)^{1/p}\qquad(m\geq2).
\]
Thus $c_m,d_m>0$ and $c_m^p+d_m^p=1$.  Define the normal Jordan
$*$-monomorphism
\[
        J((x_m))=\bigl((x_m,\theta_m(x_m))\bigr)_{m\geq2}
        \quad (x=(x_m)\in M).
\]
Let
\[
        b=\bigl((c_m1_m,d_m1_m)\bigr)_{m\geq2}\in N_+,\qquad
        w=J(1)=1_N.
\]
Since all $c_m$ and $d_m$ are strictly positive, $s(b)=1_N=J(1)=w^*w$, and
$b$ commutes with $J(M)$.  Now set
\[
        T(x)=wbJ(x)
        =\bigl((c_mx_m,d_mx_m^T)\bigr)_{m\geq2}.
\]
For $x=(x_m)\in\Lp(M)$, transposition preserves the Schatten $p$-norm, and
therefore
\[
\begin{aligned}
 \|T(x)\|_{\Lp(N)}^p
 &=\sum_{m=2}^{\infty}2^{-m}
        \bigl(c_m^p\|x_m\|_{\Sp_m}^p+d_m^p\|x_m^T\|_{\Sp_m}^p\bigr) \\
 &=\sum_{m=2}^{\infty}2^{-m}(c_m^p+d_m^p)\|x_m\|_{\Sp_m}^p
  =\|x\|_{\Lp(M)}^p .
\end{aligned}
\]
Thus $T$ is an isometry.  The displayed formula also shows that
$\tau_N(b^pJ(y))=\tau_M(y)$ for $y\in M_+$, so this is the Yeadon
factorization in the convention recalled in \cite[Section 2]{Arhancet2020}.

We next prove complete boundedness of $T$.  The source paper records the
standard formula
\[
        \|\theta_m:\Sp_m\to\Sp_m\|_{\cb}=m^{\alpha}
        \qquad\text{\cite[p. 7]{Arhancet2020}}.
\]
Fix a matrix level $k$ and an element $u$ in the $k$th amplification of the
$m$th Schatten block.  At this level the $m$th block of $T$ satisfies
\[
 \|(c_mu,d_m\theta_m^{(k)}u)\|^p
 \leq
 \bigl(c_m^p+d_m^pm^{\alpha p}\bigr)\|u\|^p
 =
 \bigl(2-m^{-\alpha p}\bigr)\|u\|^p
 \leq 2\|u\|^p .
\]
The matrix amplification of the direct product $L^p$-space is the corresponding
weighted $\ell_p$-sum of the amplified blocks.  Summing the previous estimate
over $m$ gives
\[
        \|T\|_{\cb,\Lp(M)\to\Lp(N)}\leq 2^{1/p}.
\]

On the other hand $J$ is not completely bounded.  Indeed, the projection from
$N$ onto the second coordinate of the $m$th summand is completely contractive,
and its composition with $J$ restricted to the $m$th copy of $M_m$ is exactly
$\theta_m:M_m\to M_m$.  The C$^*$-level cb norm of transposition on $M_m$ is
$m$.  Hence $\|J\|_{\cb}\geq m$ for every $m\geq2$, which is impossible for a
completely bounded map.

It remains to check that the structural conclusion of Conjecture 5.1 also
fails.  For each $m\geq2$, the restriction
\[
        T_m:\Sp_m\longrightarrow \Sp_m\oplus_p\Sp_m,\qquad
        T_m(x)=(c_mx,d_mx^T),
\]
is not a complete isometry.  To see this, choose almost extremal vectors in
matrix amplifications for the transpose norm.  Letting the approximation
error tend to zero gives
\[
        \|T_m\|_{\cb}
        \geq
        \bigl(c_m^p+d_m^pm^{\alpha p}\bigr)^{1/p}
        =
        \bigl(2-m^{-\alpha p}\bigr)^{1/p}>1.
\]
Since $T_m$ is an isometry at scalar level, this strict amplification growth
rules out complete isometry.

Any decomposition $M=M_1\oplus M_2$ is given by a central projection, hence by
a partition of the matrix factors $M_m$.  If the $m$th factor, $m\geq2$, lay in
$M_1$, then $T|_{\Lp(M_1)}$ being a complete isometry would force $T_m$ to be a
complete isometry, a contradiction.  Thus every factor $M_m$, $m\geq2$, must
lie in $M_2$.  But then $M_2$ contains full matrix algebras of arbitrarily
large dimension, so it is not $n$-minimal for any finite $n$; equivalently,
it is not $n$-subhomogeneous for any finite $n$, using the equivalence quoted
in \cite[Section 2]{Arhancet2020}.  No decomposition as in Conjecture 5.1
exists.
\end{proof}

\section*{Verification notes}

The coefficient choice was checked by
\texttt{code/check\_block\_estimates.py} for representative values of $p$ and
$2\leq m\leq80$.  This is only a sanity check: the proof above uses the exact
identities
\[
        c_m^p+d_m^p=1,\qquad d_mm^\alpha=1,\qquad
        c_m^p+d_m^pm^{\alpha p}=2-m^{-\alpha p}>1.
\]
The source crop was generated from the local PDF by
\texttt{code/make\_open\_problem\_crop.py}.

\section*{Novelty and limitations}

The local lightweight indexes were searched:
\begin{quote}
\ttfamily
registry\_index.tsv\\
solutions/index.tsv\\
attempts/index.tsv\\
proof\_gaps/index.tsv
\end{quote}
The search terms were \texttt{2007.06999}, the paper title,
\texttt{completely bounded isometries}, \texttt{Yeadon factorization},
\texttt{normal Jordan}, and \texttt{n-minimal}.  These searches found no exact
duplicate and no existing answer packet for Conjectures 5.1--5.2.

External web searches on 2026-07-18 for the arXiv id and exact phrases around
``If $T$ is completely bounded then the normal Jordan homomorphism $J$ is
completely bounded'', ``completely bounded isometries Yeadon factorization'',
and ``normal Jordan completely bounded $L^p(M)$ isometry'' found the source
paper and related citations, but no later resolution of these conjectures.
Novelty therefore remains candidate-level pending human or bibliographic
review, but no known answer was found in the bounded search.

The main human-review point is the standard operator-space bookkeeping for
matrix amplifications of direct sums of noncommutative $L^p$ spaces.  The
counterexample otherwise uses only finite matrix blocks, positive scalar
weights, and the transpose cb-norm formula explicitly recalled in the source.

\begin{thebibliography}{9}

\bibitem{Arhancet2020}
C. Arhancet,
\emph{A characterization of completely bounded normal Jordan
$*$-homomorphisms on von Neumann algebras},
arXiv:2007.06999, 2020.

\bibitem{Yeadon1977}
F. J. Yeadon,
\emph{Isometries of non-commutative $L^p$-spaces},
Math. Proc. Cambridge Philos. Soc. 90 (1981), 41--50.

\end{thebibliography}

\end{document}
